\section{数据描述与任务定义}
数据集由 30 名受试者佩戴腰部智能手机完成六类活动，采样频率为 50 Hz，采用长度为 2.56 s、含 128 个采样点且相邻窗口重叠 50\% 的窗口构造样本\cite{uci,anguita2013}。本文沿用官方按受试者划分的训练集与测试集：训练集含 21 人、7352 个窗口，测试集含另外 9 人、2947 个窗口。检查结果表明，两集合的受试者交集为空，训练特征均为有限数值；各类别数量见表\ref{tab:data}。
\begin{table}[htbp]\centering\small
\caption{官方训练集与测试集的类别分布}\label{tab:data}
\begin{tabular}{lrrr}\toprule
类别 & 训练窗口数 & 测试窗口数 & 总数\\\midrule
步行 & 1226 & 496 & 1722\\
上楼 & 1073 & 471 & 1544\\
下楼 & 986 & 420 & 1406\\
坐 & 1286 & 491 & 1777\\
站立 & 1374 & 532 & 1906\\
躺卧 & 1407 & 537 & 1944\\
合计 & 7352 & 2947 & 10299\\
\bottomrule\end{tabular}
\end{table}


记第 $i$ 个窗口为 $(\boldsymbol{x}_i,\boldsymbol{S}_i,y_i,g_i)$，其中 $\boldsymbol{x}_i\in\mathbb{R}^{561}$ 为人工特征，$\boldsymbol{S}_i\in\mathbb{R}^{9\times128}$ 为时序信号，$y_i\in\{1,\ldots,6\}$ 为活动标签，$g_i$ 为受试者编号。九个通道依次为三轴身体加速度、三轴角速度和三轴总加速度。两种输入按同一窗口的行号对齐，分类目标是学习从输入到六类活动的映射。

人工特征及预处理后的信号窗口均由数据集发布者提供，本文没有重新执行连续信号去噪、切窗或人工特征提取。神经网络实现将标签减一编码为 $0$ 至 $5$，报告时恢复原活动名称。两条分支使用同一窗口的不同表示，因而称为表示或特征融合；它们不构成独立采集的两种传感器模态。
